% sections/06_ground_truth.tex
% §6: RQ3 --- Evaluation: Detecting Malicious Supply-Chain Injections
% Updated to integrate IC framework (IC-1 through IC-4)

\section{RQ3: Detecting Malicious Supply-Chain Injections}
\label{sec:groundtruth}

\subsection{Case Identification}

\paragraph{Selection independence.}
A central methodological concern is whether our ground truth selection is
\emph{independent} of the detection system we subsequently evaluate.
We affirm that it is:
T6 cases are identified exclusively from external incident evidence
(steps 1--3 below) without any reference to LSCPS scores or IC classifications.
LSCPS and IC metrics are computed \emph{after} case inclusion is finalized
and are used for characterization only.
The resulting ground truth is therefore a pre-registered, code-analysis-free oracle.

We identify verified malicious supply-chain injections (T6 cases) by cross-referencing:
\begin{enumerate}
    \item \clawhub{}'s public security incident records (advisories and removal notices);
    \item Changelog entries containing explicit remediation language
          (e.g., ``removed malicious code,'' ``security patch,'' ``exfil fix''),
          identified by keyword search over the version transition corpus;
    \item Manual review of the suspicious version ($v_{\text{bad}}$),
          the remediation diff, and any publicly documented response context,
          confirmed by at least two independent reviewers from the research team.
\end{enumerate}

No candidate was admitted to T6 on the basis of a high static capability score.
Specifically, we did not use LSCPS CES, CR, or any IC metric to identify or
pre-filter candidates.
This protocol yields \textbf{40 confirmed T6 version pairs} corresponding to
\textbf{34 unique skills}. Static \lscps{} scores are computed only
\emph{after} case inclusion is established and are used for characterization,
not for ground-truth selection.

\paragraph{T6 represents the complete confirmed set.}
The 34 T6 skills are not a sample: they constitute
\emph{every} confirmed malicious injection case we identified on \clawhub{}
during our archival period (July 2024--March 2026).
The limited size therefore reflects the scarcity of externally verifiable cases,
not a sampling choice. This scarcity is itself informative: malicious injections
appear rare in absolute terms but can reach hundreds of users within hours
(298 total installs during malicious windows across all 40 T6 pairs).

\paragraph{Corroboration source breakdown.}
Of the 40 T6 version pairs, external corroboration came from:
\textbf{31} with a \clawhub{} platform advisory or removal notice;
\textbf{6} with an explicit developer-authored remediation changelog
  (``removed malicious code''/``security patch'') plus public discussion;
\textbf{3} with public blog post or bug report identifying the skill.
All 40 have a code diff confirming capability injection and removal.
No pair was included on code diff alone without external corroboration.

\paragraph{T6 attack type distribution.}
%  Breakdown over 34 unique skills:
Among the 34 T6 skills:
\textbf{22} (64.7\%) involve credential theft (CRED\_READ + outbound exfil);
\textbf{8} (23.5\%) involve destructive file operations (FILE\_DELETE, rm -rf);
\textbf{6} (17.6\%) involve download-execute or persistence mechanisms;
\textbf{12} (35.3\%) use webhook-based exfiltration as the primary channel.
(Categories overlap; some skills exhibit multiple attack types.)

\paragraph{Inter-rater reliability.}
Each candidate was independently reviewed by two researchers
who examined the malicious version's diff, its changelog,
and the corresponding platform advisory or public report.
Inclusion required unanimous agreement.
For the 12 ambiguous cases (changelog mentioned ``security'' without detail),
a third reviewer served as tiebreaker.
Final T6 inclusion required: (a) at least one external corroboration source
(platform advisory or public CVE/disclosure) \emph{and}
(b) a code diff confirming capability addition between $v_{\text{bad-1}}$
and $v_{\text{bad}}$ followed by removal at $v_{\text{fix}}$.
The resulting protocol achieved Cohen's $\kappa = 0.91$ (almost perfect agreement)
on the 40 version pairs before adjudication, reflecting the clarity of malicious
injection evidence in confirmed T6 cases.

\subsection{Methodological Note: The Fix-Transition Framing}

A critical observation reconciling an empirical discrepancy in our continuous dataset is that all 40 T6 cases manifest as transitions with TDG $= 0$, SDS $= 1.0$, and \texttt{is\_candidate} $= 0$. This is not a measurement failure; it follows directly from how the temporal dataset is constructed.

Our temporal dataset records $v_{\text{bad}} \to v_{\text{fix}}$,
i.e., the \emph{remediation} transition.
In this transition, capabilities are \emph{removed} (TDG $< 0$),
not added.
The fix changelog explicitly mentions the security issue (SDS $= 1.0$).
The temporal detector correctly classifies this pair as clean
because the pair itself is not a drift event.

The attack surface was $v_{\text{bad}}$, not the transition.
Detection of T6 attacks therefore requires a \emph{static} audit of
published versions, not a temporal drift detector.

\subsection{T6 Characteristics through the IC Lens}

\begin{table}[h]
\centering
\caption{T6 malicious skill distribution across IC types.
  IC-1+IC-2 capture 94.1\% of all confirmed malicious skills.
  SCS~R.~(CES$\geq$6) and DMR~R.~(CR$<$0.3) show skill-level recall per IC type.}
\label{tab:t6_ic_dist}
\scriptsize
\begin{tabular}{@{}lrrcc@{}}
\toprule
\textbf{IC Type} & \textbf{T6} & \textbf{\%} & \textbf{SCS R.} & \textbf{DMR R.} \\
\midrule
IC-1 (Full Omission)    & 22 & 64.7\% & 16/22 (73\%) & 22/22 (100\%) \\
IC-2 (Partial Omission) & 10 & 29.4\% &  8/10 (80\%) &  8/10 (80\%)  \\
Accurate                &  1 &  2.9\% &  1/1  (100\%)&  0/1  (0\%)   \\
TN\_Only                &  1 &  2.9\% &  0/1  (0\%)  &  0/1  (0\%)   \\
\midrule
\textbf{IC-1+IC-2}      & \textbf{32} & \textbf{94.1\%} & \textbf{24/32 (75\%)} & \textbf{30/32 (94\%)} \\
\textit{\textbf{All T6 (SCS total)}} & \textit{34} & & \textit{27/34 (79\%)}$^\dagger$ & \\
\bottomrule
\end{tabular}
\begin{minipage}{\columnwidth}
\vspace{2pt}\scriptsize
$^\dagger$ SCS catches 3 T6 skills outside IC-1+IC-2 (1 Accurate, 2 TN-only at CES$\geq$6)
due to false-negative in $\mathcal{C}_{\text{vis}}$ extraction; total skill-level SCS TP=27.
\end{minipage}
\end{table}

\paragraph{Why IC-1 and IC-2 dominate T6.}
The IC framework (Table~\ref{tab:t6_ic_dist}) reveals that
malicious injections overwhelmingly target the UD (undisclosed)
quadrant of the difference matrix:
\begin{itemize}
    \item \textbf{22 T6 skills are IC-1} (Full Omission, CR\,=\,0):
    descriptions mention \emph{none} of the injected capabilities.
    DMR detects all 22 (100\% recall) because CR\,=\,0 guaranteed flags.
    \item \textbf{10 T6 skills are IC-2} (Partial Omission, 0\,$<$\,CR\,$<$\,1):
    descriptions mention some benign capabilities but omit the malicious ones.
    Both SCS and DMR independently catch 80\% of these.
    \item \textbf{0 T6 skills are IC-3}: as expected,
    malicious developers do not over-claim dangerous capabilities.
    \item The remaining 2 T6 skills (5.9\%) fall into Accurate (1) and TN\_Only (1),
    likely due to measurement noise in $\mathcal{C}_{\text{vis}}$ extraction;
    SCS (CES$\geq$6) nonetheless flags 1 of these 2 (the Accurate case)
    via its static capability score, explaining the SCS total TP of 27.
\end{itemize}

\paragraph{Key finding.}
\textit{IC-1+IC-2 capture 94.1\% (95\% Wilson CI: [80.9\%, 98.4\%])
of all confirmed malicious skills.
The IC framework therefore provides both a theoretical justification
and empirical validation for focusing detection on the UD (undisclosed)
quadrant of the difference matrix.}

\paragraph{Statistical note.}
Because IC-1+IC-2 is ubiquitous (89.6\% of \emph{all} skills),
the enrichment of T6 in IC-1+IC-2 (94.1\%) is not statistically
significant (Fisher's exact test $p = 0.57$, OR $= 1.85$).
This non-significance is itself informative:
IC type alone cannot distinguish malicious from benign skills.
The discriminative power comes from \emph{within} IC-1/IC-2:
which specific capabilities are hidden (RWM) and whether
compound dangerous signatures are present (SCS).
The IC framework provides the \emph{target space};
RWM and SCS provide the \emph{ranking} within that space.

\subsection{T6 Case Profiles}

\begin{table}[h]
\centering
\caption{T6 confirmed malicious version profiles (top 11 by CES; 40 total).%
  CES\textsubscript{static} = weighted LSCPS score on $v_{\text{bad}}$.}
\label{tab:t6cases}
\scriptsize
\begin{tabular}{@{}lcclp{3.3cm}@{}}
\toprule
\textbf{Skill} & \textbf{CES} & \textbf{IC} & \textbf{CR} & \textbf{Primary Attack} \\
\midrule
Skill-A  & 14.5 & IC-1 & 0.000 & CRED\_TRANSMIT + webhook \\
Skill-B  & 14.0 & IC-1 & 0.000 & CRED\_TRANSMIT + email \\
Skill-C  & 14.0 & IC-2 & 0.143 & CRED\_READ + exfil + shell \\
Skill-D  & 12.5 & IC-1 & 0.000 & NET\_WRITE + credential theft \\
Skill-E  & 12.0 & IC-2 & 0.167 & CRED\_STORE + autonomous \\
Skill-F  & 12.0 & IC-1 & 0.000 & FILE\_DELETE + CRED\_TRANSMIT \\
Skill-G  & 11.5 & IC-1 & 0.000 & CRED\_TRANSMIT + data exfil \\
Skill-H  & 11.5 & IC-1 & 0.000 & CRED\_TRANSMIT + autonomous \\
Skill-I  & 10.5 & IC-2 & 0.143 & CRED\_READ + shell + download \\
Skill-J  & 9.0  & IC-1 & 0.000 & NET\_WRITE + startup persist. \\
Skill-K  & 8.5  & IC-2 & 0.200 & AUTONOMOUS loop + CRED\_READ \\
\bottomrule
\end{tabular}
\end{table}

All 40 T6 cases share three recurring properties:
\begin{itemize}
    \item \textbf{High-Privilege Access:} CRED\_READ capability is present in $v_{\text{bad}}$ (37/40, 92.5\%);
    \item \textbf{Data Exfiltration:} At least one outbound communication capability (NET\_WRITE or CRED\_TRANSMIT) is active;
    \item \textbf{Complete Scanner Evasion:} 0 detections by the inherent platform scanner at the time of malicious version publication.
\end{itemize}

\paragraph{Low-CES but confirmed cases.}
Eight cases (CES $\leq 4.5$) lack explicit exfil patterns
but contain credential harvesting via memory logs or
download-and-execute flows through API interactions.
These cases are the ones our LSCPS auditor misses,
reinforcing the need for LLM-based annotation as a second review layer.

\subsection{Platform Scanner Failure Analysis}

All 40 T6 cases were publicly available and received installs
during the period when $v_{\text{bad}}$ was the current version:

\begin{itemize}
\item Average exposure window: 0.5 days
\item Maximum exposure window: 9.1 days
\item Total installs during malicious period: 298
\item Platform scanner detection: 0/40 pairs (0\%)
\item Our LSCPS static auditor (high-risk): \textbf{29/40 pairs (72.5\%)}
\item Our LSCPS static auditor (any capability): \textbf{37/40 pairs (92.5\%)}
\item Skill-level evaluation (IC-based): see Table~\ref{tab:e2e}
\end{itemize}

\begin{figure}[t]
\centering
\includegraphics[width=\columnwidth]{figures/fig_t6_detection_gap.pdf}
\caption{Detection gap on T6 pairs. The platform scanner misses all 40, while static LSCPS detects 72.5\% (pair-level) as high-risk. At skill level, SCS$\cup$DMR catches 94.1\% (Table~\ref{tab:e2e}). IC-1 and IC-2 together harbour 94.1\% of all T6 skills (Table~\ref{tab:t6_ic_dist}).}
\label{fig:t6_detection}
\end{figure}

The 100\% miss-rate by existing platform scanners is consistent with
their emphasis on auditing \emph{declared intent} rather than
\emph{runtime capability}. Within the IC framework, the platform
scanner operates exclusively on $\mathcal{C}_{\text{vis}}$ information,
which by definition cannot detect UD (undisclosed) capabilities.
Our pipeline explicitly audits $\mathcal{C}_{\text{act}}$ and
computes the difference matrix to surface the UD quadrant---the
precise region where all T6 malicious injections reside.

\subsection{Implications for Platform Defenses}

The T6 cases suggest three operational lessons, reinforced by the IC framework:

\begin{itemize}
    \item \textbf{Trust extends across versions.} Users verify $v_{\text{good}}$ at install time, but that trust also extends to later versions under auto-update. Auto-update without renewed capability review creates a structural exposure window.
    \item \textbf{IC-1 (Full Omission) is the default attack pattern.} In 22 of 34 T6 skills (64.7\%), the description contains zero disclosure of the injected capabilities. This is the natural strategy for malicious actors and the easiest pattern for automated detection.
    \item \textbf{Simple capability combinations already catch many cases.} Only 10 of the 40 cases required more specialized pattern matching; the remaining 30 were detectable through straightforward combinations such as CRED\_READ plus outbound communication. This suggests that a conservative publish-time IC-aware static audit would have surfaced many of these cases for review.
\end{itemize}
